\chapter{Joint Pain and Swelling}

\section*{Definition and Overview}
Joint pain with swelling is a combination of symptoms that indicates something is actively wrong inside the joint. The swelling may be due to fluid, inflammation of the joint lining, or injury to cartilage and ligaments, and its presence changes the urgency of assessment. The most important causes are inflammatory arthritis, gout, injury, and infection. A hot, swollen, red joint with fever is septic arthritis until proven otherwise and requires urgent treatment, because infection can destroy a joint within days. Other causes, such as rheumatoid arthritis and gout, respond well to early treatment, so persistent joint swelling should always be assessed.

\section*{Epidemiology}
Swelling accompanies joint pain in a minority but important proportion of joint problems. Gout affects around 7 per cent of Australian men and causes recurrent flares of a swollen joint, most often the big toe. Rheumatoid arthritis, affecting around 1 per cent of people, typically presents with swelling of the hands and wrists. Septic arthritis is rare but serious, occurring most often in children, older people, and people with weakened immunity or pre-existing joint disease. Traumatic swelling follows injury and is very common in active people.

\section*{Aetiology and Pathophysiology}
Joint swelling occurs when fluid or inflammation accumulates in the joint. In inflammatory arthritis, the immune system attacks the synovium, which thickens and produces excess fluid, with swelling, warmth, and stiffness. In gout, uric acid crystals in the joint trigger intense acute inflammation with marked swelling. In septic arthritis, bacteria multiply in the joint fluid, causing rapid destruction of cartilage. In injury, the joint swells from bleeding or fluid following ligament, cartilage, or bone damage. The location of the swelling, the number of joints affected, and systemic features such as fever help identify the cause and guide urgent versus routine care.

\section*{Clinical Features}
\begin{itemize}
    \item Visible swelling of the joint, with a feeling of fullness or tightness
    \item Pain, which may be constant or with movement
    \item Warmth and redness over the joint
    \item Reduced range of movement and stiffness
    \item A hot, red, swollen joint with fever, which may indicate septic arthritis
    \item Symmetrical swelling of small joints, suggesting rheumatoid arthritis
    \item Sudden severe swelling of a single joint, suggesting gout
    \item Swelling after injury, with bruising and tenderness
\end{itemize}

\section*{Differential Diagnosis}
The cause of a swollen, painful joint must be distinguished.
\begin{itemize}
    \item \textbf{Septic arthritis:} a hot, swollen joint with fever, an emergency
    \item \textbf{Gout and pseudogout:} sudden flares, often in the big toe or knee
    \item \textbf{Rheumatoid arthritis:} symmetrical swelling with morning stiffness
    \item \textbf{Psoriatic and reactive arthritis:} other inflammatory forms
    \item \textbf{Injury:} ligament, cartilage, and fracture-related swelling
    \item \textbf{Osteoarthritis:} generally mild swelling without significant inflammation
    \item \textbf{Bursitis and tendonitis:} swelling around, not within, the joint
    \item \textbf{Haemarthrosis:} bleeding into the joint after injury or in clotting disorders
\end{itemize}

\section*{When to Seek Medical Care}
Joint pain with swelling always warrants assessment. Urgent care is needed for a hot, swollen, red joint, especially with fever, for swelling after significant injury, and for swelling that is sudden and severe. Persistent swelling in multiple joints, or with morning stiffness and fatigue, warrants early specialist review, as early treatment of inflammatory arthritis prevents damage.

\textbf{Call 000 (or local emergency services) for:}
\begin{itemize}
    \item A hot, swollen, red joint with fever, which may indicate septic arthritis
    \item Sudden severe joint swelling after injury, with deformity or inability to use the joint
    \item Sudden breathlessness or chest pain in a person with inflammatory arthritis
    \item Collapse, confusion, or severe illness
\end{itemize}

\section*{Diagnosis and Investigations}
\begin{itemize}
    \item \textbf{History and examination:} the joints involved, onset, and systemic features
    \item \textbf{Joint aspiration:} examination of the fluid for crystals, infection, and inflammation
    \item \textbf{Blood tests:} inflammatory markers, rheumatoid factor, anti-CCP, and uric acid
    \item \textbf{X-rays:} for bone and joint changes
    \item \textbf{Ultrasound and MRI:} for inflammation and soft tissue injury
    \item \textbf{Microbiology:} culture of joint fluid when infection is suspected
    \item \textbf{Referral:} urgent to rheumatology or orthopaedics as indicated
\end{itemize}

\section*{Management}
\begin{itemize}
    \item \textbf{Septic arthritis:} urgent intravenous antibiotics and joint drainage
    \item \textbf{Gout:} anti-inflammatory treatment and urate-lowering therapy
    \item \textbf{Rheumatoid arthritis:} disease-modifying medicines started early
    \item \textbf{Injury:} rest, ice, compression, elevation, and assessment for ligament or fracture
    \item \textbf{Pain relief:} anti-inflammatories and simple analgesics
    \item \textbf{Joint injection:} corticosteroids after infection is excluded
    \item \textbf{Physiotherapy:} rehabilitation and joint protection
    \item \textbf{Follow-up:} monitoring of the response and ongoing management
\end{itemize}

\section*{Complications}
\begin{itemize}
    \item Destruction of the joint in septic arthritis if treatment is delayed
    \item Joint damage and deformity in untreated inflammatory arthritis
    \item Recurrent gout flares with joint damage over time
    \item Chronic pain and loss of function
    \item Complications of injury, including instability and arthritis later
    \item Side effects of treatment
\end{itemize}

\section*{Prognosis and Outcomes}
The outlook depends on the cause and the speed of treatment. Septic arthritis has an excellent outcome when treated urgently, with joint damage avoided. Gout resolves with treatment, and inflammatory arthritis is well controlled with modern medicines, especially when started early. Injury-related swelling settles with rehabilitation. Because joint swelling is a sign of active disease in the joint, it should never be ignored, and early assessment gives the best outcomes.

\section*{Prevention and Self-Care}
\begin{itemize}
    \item Do not ignore a swollen joint: seek assessment promptly
    \item If a joint is hot, red, and painful, seek urgent care
    \item Protect joints from injury with technique and equipment
    \item Manage gout risk factors: diet, alcohol, and urate control
    \item Take prescribed medicines for inflammatory arthritis
    \item Use ice, elevation, and rest for injury-related swelling
    \item Attend rehabilitation and follow-up
    \item Report new or changing symptoms early
\end{itemize}

\section*{Sources and Further Reading}
The following reputable sources provide further information for healthcare students and patients seeking reliable, current advice about joint pain and swelling.
\begin{itemize}
    \item Arthritis Australia. \textit{Swollen joints and arthritis.} https://arthritisaustralia.com.au/
    \item Healthdirect Australia. \textit{Joint swelling.} https://www.healthdirect.gov.au/
    \item Australian Rheumatology Association. \textit{Inflammatory arthritis.} https://www.rheumatology.org.au/
    \item Better Health Channel. \textit{Arthritis and joint problems.} https://www.betterhealth.vic.gov.au/
    \item NHS. \textit{Joint pain and swelling.} https://www.nhs.uk/
    \item Mayo Clinic. \textit{Swollen joints.} https://www.mayoclinic.org/
\end{itemize}
